\documentclass[11pt]{article}
\usepackage[margin=1in]{geometry}
\usepackage{amsmath,amssymb,amsthm,mathtools}
\usepackage{booktabs}
\usepackage{microtype}
\usepackage[hidelinks]{hyperref}
\usepackage{enumitem}

\newtheorem{theorem}{Theorem}[section]
\newtheorem{lemma}[theorem]{Lemma}
\newtheorem{proposition}[theorem]{Proposition}
\newtheorem{corollary}[theorem]{Corollary}
\theoremstyle{remark}
\newtheorem{remark}[theorem]{Remark}

\newcommand{\per}{\operatorname{per}}
\newcommand{\R}{\mathbb{R}}
\newcommand{\one}{\mathbf{1}}
\newcommand{\Om}{\Omega}

\title{Dittert's Conjecture in Dimensions 14, 15, and 16\\
via a Two-Zero Boundary Gap}
\author{Pedro Paulo Marques do Nascimento}
\date{23 July 2026}

\hypersetup{
  pdftitle={Dittert's Conjecture in Dimensions 14, 15, and 16 via a Two-Zero Boundary Gap},
  pdfauthor={Pedro Paulo Marques do Nascimento},
  pdfsubject={Exact working proof and reproducible verification},
  pdfkeywords={Dittert conjecture, permanent, exact verification, boundary gap}
}

\begin{document}
\maketitle

\begin{abstract}
Let $K_n$ be the set of nonnegative $n\times n$ matrices whose entries sum to $n$, and let
\[
 \Phi(A)=\prod_{i=1}^n r_i+\prod_{j=1}^n c_j-\per(A),
\]
where $r_i$ and $c_j$ are the row and column sums.  We prove, by a short exact
computer-assisted calculation, that Dittert's conjecture holds for
$n=14,15,16$:
\[
 \Phi(A)\le 2-\frac{n!}{n^n},
\]
with equality only at the matrix all of whose entries are $1/n$.
The proof combines the joint-deficit scaling lemma with two facts about the
boundary: a maximizing matrix must contain two zeroes in distinct rows and
columns, and the minimum permanent on that two-zero face of the Birkhoff
polytope reduces to a one-variable polynomial.  Exact rational certificates
for the three polynomial lower bounds and the final scalar comparisons are
included with this note.
\end{abstract}

\section{Setup and known results}

Let
\[
 K_n=\left\{A=(a_{ij})\in\R_{\ge0}^{n\times n}:\sum_{i,j}a_{ij}=n\right\},
 \qquad
 \gamma_n=\frac{n!}{n^n},
\]
and put $U_n=n^{-1}\one\one^{\mathsf T}$.  Dittert's conjecture asserts that
\begin{equation}\label{eq:dittert}
 \Phi(A)\le \Phi(U_n)=2-\gamma_n,
\end{equation}
with equality only at $U_n$.

We use the following established results.
\begin{enumerate}[label=(\roman*)]
 \item Every positive global maximizer of $\Phi$ on $K_n$ is $U_n$ \cite{Hwang1986}.
 \item A partly decomposable matrix cannot be a global maximizer of $\Phi$, and
 the complete zero set of a global maximizer cannot be one rectangular block
 \cite[Theorems 3.2 and 3.7]{CheonWanless2012}.
 \item A nonnegative matrix $S$ is doubly superstochastic if and only if
 \begin{equation}\label{eq:li}
   s(S[I,J])\ge |I|+|J|-n
 \end{equation}
 for every pair of row and column sets $I,J$; here $s$ denotes the sum of the
 indicated submatrix \cite[Lemma 2.2]{CheonWanless2012}.
\end{enumerate}

\section{A boundary maximizer has two independent zeroes}

\begin{lemma}\label{lem:zeros}
Let $Z$ be a nonempty set of positions in an $n\times n$ array.  If $Z$ does
not contain two positions in distinct rows and distinct columns, then $Z$ is
contained in one row or in one column.
\end{lemma}

\begin{proof}
Choose $(i_0,j_0)\in Z$.  If every element of $Z$ lies in row $i_0$, there is
nothing to prove.  Otherwise choose $(i_1,j_0)\in Z$ with $i_1\ne i_0$; the
column must be $j_0$ by the hypothesis.  If $(i,j)\in Z$ had $j\ne j_0$, then
it would have to share a row with both $(i_0,j_0)$ and $(i_1,j_0)$, which is
impossible.  Hence every element of $Z$ lies in column $j_0$.
\end{proof}

\begin{corollary}\label{cor:twozeros}
Every boundary global maximizer of $\Phi$ contains two zero entries in
distinct rows and distinct columns.
\end{corollary}

\begin{proof}
By Lemma~\ref{lem:zeros}, otherwise all zeroes lie in one row or one column.
If that row or column is entirely zero, the matrix is partly decomposable.  If
not, the complete zero set is one proper rectangular block.  Both alternatives
are excluded by the results quoted in Section~1.
\end{proof}

\section{Joint-deficit scaling}

The following is the joint-deficit scaling observation used independently for
the case $n=16$ in \cite{Kafidov2026}.  We reproduce the proof because the
shared deficit budget is central to the extension to $n=14,15$.

\begin{lemma}[Joint-deficit scaling]\label{lem:scaling}
Suppose $A\in K_n$ satisfies
\[
 \Phi(A)\ge 2-\gamma_n,
 \qquad
 \per(A)=\gamma_n-\delta,
 \qquad 0\le\delta<1.
\]
Set
\[
 t=\sqrt{\frac{n\delta}{1-\delta}}.
\]
If $t<1$, then $(1-t)^{-1}A$ is doubly superstochastic.
\end{lemma}

\begin{proof}
Write
\[
 R=\prod_i r_i=1-\rho,
 \qquad
 C=\prod_j c_j=1-\sigma.
\]
AM--GM gives $R,C\le1$.  The assumed lower bound on $\Phi$ gives
\begin{equation}\label{eq:budget}
 \rho+\sigma\le\delta.
\end{equation}

Let $I$ be a set of $k$ rows and write
$\sum_{i\in I}r_i=k+\varepsilon$.  For $k=0$ or $k=n$ we have
$\varepsilon=0$, so the required deviation bound is immediate.  We may
therefore assume $1\le k\le n-1$.  Applying AM--GM inside $I$ and its
complement gives
\[
 R\le\left(1+\frac{\varepsilon}{k}\right)^k
       \left(1-\frac{\varepsilon}{n-k}\right)^{n-k}.
\]
With $p=k/n$ and $q=(k+\varepsilon)/n$, the logarithm of the right-hand side
is $-nD(p\|q)$.  Binary Pinsker, $D(p\|q)\ge2(p-q)^2$, therefore yields
\[
 |\varepsilon|
 \le \sqrt{-\frac n2\log(1-\rho)}
 \le \sqrt{\frac{n\rho}{2(1-\delta)}}.
\]
The same estimate holds for every column subset, with $\sigma$ in place of
$\rho$.

Now take row and column sets $I,J$ and put
$p_0=|I|+|J|-n>0$.  By inclusion--exclusion and nonnegativity,
\begin{align*}
 s(A[I,J])
 &\ge p_0
 -\sqrt{\frac{n\rho}{2(1-\delta)}}
 -\sqrt{\frac{n\sigma}{2(1-\delta)}}\\
 &\ge p_0-\sqrt{\frac{n(\rho+\sigma)}{1-\delta}}
 \ge p_0-t,
\end{align*}
where we used \eqref{eq:budget}.  Since $p_0$ is a positive integer,
$p_0-t\ge(1-t)p_0$.  Thus $(1-t)^{-1}A$ satisfies \eqref{eq:li}.  Pairs with
$p_0\le0$ satisfy it automatically.
\end{proof}

\section{The permanent gap on the two-zero face}

Let $\Omega_n^{(2)}$ be the set of doubly stochastic $n\times n$ matrices
having prescribed zeroes in two distinct rows and columns.  Permuting columns,
we may take those zeroes to be positions $(1,2)$ and $(2,1)$.

Put $m=n-2$.  In the notation of Pula--Song--Wanless this is the face
$\Omega(V_{m,2})$ \cite{PulaSongWanless2011}.  Minc's averaging theorem for
permanent minimizers with repeated support rows and columns implies that a
minimizer on this face may be chosen in the block form
\cite[Theorem~1]{Minc1984}; see also \cite{PulaSongWanless2011}
\begin{equation}\label{eq:Ygeneral}
 Y=
 \begin{pmatrix}
 x_1&0&a_1\one_m^{\mathsf T}\\
 0&x_2&a_2\one_m^{\mathsf T}\\
 a_1\one_m&a_2\one_m&xJ_m
 \end{pmatrix},
 \qquad
 x_i=1-ma_i,
 \qquad
 a_1+a_2+mx=1.
\end{equation}
Here $0\le a_i\le1/m$.  Only the averaging step of Minc's theorem is used to
obtain \eqref{eq:Ygeneral}.  Section~2 of Pula--Song--Wanless assumes both block
parameters exceed $2$, so its equalization theorem is not invoked for the
present $V_{m,2}$ face.  Although Minc already resolved this face, the following
short calculation supplies the one-variable reduction directly and avoids
importing any unrecorded formula for the minimizer.

\begin{lemma}[Equalization of the exceptional parameters]\label{lem:equalize}
For $m\ge4$, among matrices of the form \eqref{eq:Ygeneral} with fixed
$s=a_1+a_2$, the permanent is minimized when $a_1=a_2=s/2$.
\end{lemma}

\begin{proof}
Put $p=a_1a_2$ and $x=(1-s)/m$.  Counting permanent terms according to whether
both, one, or neither of the two entries $x_1,x_2$ are used gives
\begin{equation}\label{eq:generalper}
 \per(Y)=m!x^{m-2}\Bigl[
 x^2x_1x_2+mx(x_1a_2^2+x_2a_1^2)+m(m-1)a_1^2a_2^2
 \Bigr].
\end{equation}
The symmetric identities
\[
 x_1x_2=1-ms+m^2p,
 \qquad
 x_1a_2^2+x_2a_1^2=s^2-(2+ms)p
\]
show that the bracket in \eqref{eq:generalper}, as a polynomial in $p$, has
derivative
\begin{equation}\label{eq:pderivative}
 2m(m-1)p+(m+1)s^2-ms-1.
\end{equation}
The constraints $0\le a_i\le1/m$ imply $0\le s\le2/m$, while
$p\le s^2/4$.  Since \eqref{eq:pderivative} increases with $p$, it is at most
\[
 D_m(s)=\frac{(m^2+m+2)s^2-2ms-2}{2}.
\]
The quadratic $D_m$ is convex, and
\[
 D_m(0)=-1,
 \qquad
 D_m(2/m)=-1+\frac2m+\frac4{m^2}<0
 \quad(m\ge4).
\]
Hence \eqref{eq:pderivative} is strictly negative throughout the feasible
region.  For fixed $s$, increasing $p$ therefore decreases the permanent, and
the largest possible value $p=s^2/4$ is attained exactly at
$a_1=a_2=s/2$.
\end{proof}

It follows that a minimizing matrix may be chosen as
\begin{equation}\label{eq:Bmx}
 B_m(x)=
 \begin{pmatrix}
 a&0&b\one_m^{\mathsf T}\\
 0&a&b\one_m^{\mathsf T}\\
 b\one_m&b\one_m&xJ_m
 \end{pmatrix},
 \quad
 a=\frac{2-m+m^2x}{2},\quad b=\frac{1-mx}{2},
\end{equation}
where
\begin{equation}\label{eq:xrange}
 \frac{m-2}{m^2}\le x\le\frac1m.
\end{equation}
A second reading of the three cases in \eqref{eq:generalper} gives
\begin{equation}\label{eq:simplePn}
 \per B_m(x)=m!x^{m-2}
 \bigl(x^2a^2+2mxab^2+m(m-1)b^4\bigr).
\end{equation}
Equivalently,
\begin{equation}\label{eq:Pn}
 P_n(x):=\per B_m(x)
 =\frac{(m!)^2}{16}\sum_{i=0}^2
 \binom2i\frac{2^i}{(m-2+i)!}
 (2-m+m^2x)^i(1-mx)^{4-2i}x^{m-2+i}.
\end{equation}

For the three dimensions, exact expansion gives
\begin{align*}
P_{14}(x)&=119750400x^{10}
 (953856x^4-289728x^3+33220x^2-1704x+33),\\
P_{15}(x)&=1556755200x^{11}
 (1513733x^4-427739x^3+45582x^2-2171x+39),\\
P_{16}(x)&=10897286400x^{12}
 (4648336x^4-1227744x^3+122200x^2-5432x+91).
\end{align*}
Their derivatives factor as
\begin{align*}
P'_{14}(x)&=718502400x^9(84x-5)
 (26496x^3-5896x^2+440x-11),\\
P'_{15}(x)&=20237817600x^{10}(195x-11)
 (8957x^3-1857x^2+129x-3),\\
P'_{16}(x)&=305124019200x^{11}(56x-3)
 (47432x^3-9204x^2+598x-13).
\end{align*}
The three cubic factors are strictly increasing on $\mathbb R$: the
discriminants of their derivatives are respectively
\[
 -847616,\qquad -71640,\qquad -1517568.
\]
The linear factors are positive throughout the intervals \eqref{eq:xrange},
and the cubic endpoint values are
\[
\begin{array}{c|cc}
 n&c_n((m-2)/m^2)&c_n(1/m)\\ \hline
14&-1/216&1/18\\
15&-2/2197&2/169\\
16&-1/343&2/49.
\end{array}
\]
Thus each $P_n$ has one global minimum on its admissible interval.  Exact
rational brackets for the critical points are
\begin{align*}
 n=14:&\quad
 \frac{70587793823672}{10^{15}}<x_*<
 \frac{70587793823673}{10^{15}},\\
 n=15:&\quad
 \frac{65989565762012}{10^{15}}<x_*<
 \frac{65989565762013}{10^{15}},\\
 n=16:&\quad
 \frac{61946716891694}{10^{15}}<x_*<
 \frac{61946716891695}{10^{15}}.
\end{align*}
Exact rational interval-Horner evaluation on these brackets gives
\begin{equation}\label{eq:Ln}
 \begin{array}{c|c|c}
 n&\displaystyle \min_{B\in\Omega_n^{(2)}}\per(B)> & L_n\\ \hline
14&7890634087/10^{15}&789/10^8\\
15&3001151594/10^{15}&3001/10^9\\
16&1139154849/10^{15}&1139/10^9.
 \end{array}
\end{equation}
The first verifier supplied with this note checks the displayed
factorizations, monotonicity, critical-point brackets, and outward rational
interval evaluations using only integer and rational arithmetic.  An
independent verifier instead applies Sturm's theorem directly to
$P_n(x)-L_n$ on the full interval \eqref{eq:xrange}.

\section{Boundary exclusion}

\begin{proposition}\label{prop:criterion}
Fix $n\ge4$.  Suppose every $B\in\Omega_n^{(2)}$ satisfies $\per(B)\ge L$ and
\begin{equation}\label{eq:eta}
 \eta:=L-\gamma_n-\frac{n^3L^2}{4(1-\gamma_n)}>0.
\end{equation}
Then a global maximizer of $\Phi$ on $K_n$ cannot lie on the boundary.
\end{proposition}

\begin{proof}
Let $A$ be a boundary global maximizer.  Since
$\Phi(A)\ge\Phi(U_n)$ and the two product terms are at most $1$, write
\[
 \per(A)=\gamma_n-\delta,
 \qquad 0\le\delta\le\gamma_n.
\]
If $\delta=0$, both product bounds are equalities, so $A$ is doubly
stochastic; the equality case in the van der Waerden theorem then gives
$A=U_n$, contrary to the assumed zero.  Hence $0<\delta\le\gamma_n$.

By Corollary~\ref{cor:twozeros}, $A$ has two independent zeroes.  Set
$t=\sqrt{n\delta/(1-\delta)}$.  For $n=14,15,16$ we have
$(n+1)\gamma_n<1$, and therefore $t<1$.  Lemma~\ref{lem:scaling} gives a
doubly stochastic matrix $B\le (1-t)^{-1}A$.  The two zeroes of $A$ are also
zeroes of $B$, so $\per(B)\ge L$.  Monotonicity and homogeneity of the
permanent give
\begin{equation}\label{eq:lowerper}
 \gamma_n-\delta=\per(A)\ge L(1-t)^n.
\end{equation}

On the other hand, Bernoulli's inequality and $\delta\le\gamma_n$ imply
\begin{align*}
L(1-t)^n-(\gamma_n-\delta)
&\ge L-\gamma_n+\delta-nL\sqrt{\frac{n\delta}{1-\delta}}\\
&\ge L-\gamma_n+x^2-nL\sqrt{\frac{n}{1-\gamma_n}}\,x\\
&\ge L-\gamma_n-\frac{n^3L^2}{4(1-\gamma_n)}
=\eta>0,
\end{align*}
where $x=\sqrt\delta$ and the final step completes the square.  This
contradicts \eqref{eq:lowerper}.
\end{proof}

For the values in \eqref{eq:Ln}, exact rational arithmetic gives
\[
\begin{array}{c|c}
 n&\eta_n\\ \hline
14&\displaystyle
\frac{1801777553768587612545302097}
{957898641732692266551875000000000000}>0\\[3mm]
15&\displaystyle
\frac{172396552001926445448721}
{24213708253338147072000000000000}>0\\[3mm]
16&\displaystyle
\frac{4164538286276424752117963756736721}
{1208924448418671074738176000000000000000000}>0.
\end{array}
\]
Numerically these margins are approximately
$1.88097\cdot10^{-9}$, $7.11979\cdot10^{-9}$, and
$3.44483\cdot10^{-9}$, respectively.

\begin{theorem}\label{thm:main}
For $n\in\{14,15,16\}$ and every $A\in K_n$,
\[
 \Phi(A)\le 2-\frac{n!}{n^n},
\]
with equality if and only if $A=U_n$.
\end{theorem}

\begin{proof}
A global maximizer exists by compactness.  Proposition~\ref{prop:criterion}
and the exact checks above exclude a boundary maximizer.  Every global
maximizer is therefore positive, and Hwang's theorem identifies it uniquely
as $U_n$.
\end{proof}

\begin{corollary}
Together with Pang's result for $n\ge17$ \cite{Pang2026}, Dittert's conjecture
holds for every $n\ge14$.
\end{corollary}

\section{Reproducibility}

The accompanying files are:
\begin{itemize}
 \item \texttt{verify\_dittert\_n14\_n16.py}: a standard-library-only exact
 verifier.  It reconstructs the two-parameter expression
 \eqref{eq:generalper}, verifies Lemma~\ref{lem:equalize}, reconstructs
 \eqref{eq:Pn}, checks all displayed factorizations, proves the critical-point
 brackets and interval lower bounds, and verifies \eqref{eq:eta} as a
 rational inequality.
 \item \texttt{audit\_dittert\_n14\_n16\_sympy.py}: an independent exact audit.
 It derives the block permanent separately, verifies the equalization step,
 checks the permanent polynomial by exact Ryser evaluation at $n+1$ rational
 points, and uses Sturm root counts to prove $P_n-L_n>0$ on each complete
 admissible interval.
\end{itemize}
Neither program uses a floating-point number in a correctness decision.
Both terminate with a certification message for all three dimensions.

\begin{thebibliography}{9}
\bibitem{CheonWanless2012}
G.-S. Cheon and I. M. Wanless,
\emph{Some results towards the Dittert conjecture on permanents},
Linear Algebra Appl. 436 (2012), 791--801.

\bibitem{Hwang1986}
S.-G. Hwang,
\emph{A note on a conjecture on permanents},
Linear Algebra Appl. 76 (1986), 31--44.

\bibitem{Kafidov2026}
B. Kafidov,
\emph{Dittert's conjecture in dimension 16 via a joint-deficit scaling lemma},
arXiv:2607.19439v1 (2026).

\bibitem{Minc1984}
H. Minc,
\emph{Minimum permanents of doubly stochastic matrices with prescribed zero
entries},
Linear and Multilinear Algebra 15 (1984), 225--243.

\bibitem{Pang2026}
Z. Pang,
\emph{Proof of Dittert's conjecture for dimensions $n\ge17$},
arXiv:2606.01531v1 (2026).

\bibitem{PulaSongWanless2011}
K. Pula, S.-Z. Song, and I. M. Wanless,
\emph{Minimum permanents on two faces of the polytope of doubly stochastic matrices},
Linear Algebra Appl. 434 (2011), 232--238.
\end{thebibliography}

\end{document}
